\documentclass[11pt]{article}
\usepackage[a4paper,margin=0.65in]{geometry}
\usepackage[T1]{fontenc}
\usepackage{amsmath,amssymb,amsfonts}
\usepackage{graphicx}
\usepackage{pdfpages}
\usepackage[english,shorthands=off]{babel}
\usepackage{fancyhdr}
\usepackage[bookmarks=false,colorlinks=false]{hyperref}
\usepackage[protrusion=true,expansion=false]{microtype}
\emergencystretch=3em
\tolerance=600
\providecommand{\modd}{\mathrm{mod}}
\providecommand{\Disc}{\mathrm{Disc}}
\providecommand{\canont}{}
\providecommand{\asterism}{*}
\providecommand{\Hessian}{H}
\providecommand{\tome}{}
\providecommand{\page}{p.}
\pagestyle{fancy}\fancyhf{}\fancyhead[L]{Pages 351--375}\fancyhead[R]{Cayley --- Collected Papers, Vol. XII}\renewcommand{\headrulewidth}{0.4pt}
\setlength{\headheight}{15pt}

\title{Collected Mathematical Papers, Vol. XII\\ \large PDF pages 351--375 (printed pp.\ 331--355)\\ Papers 845 (concl.), 846, 847 (start)}
\author{Arthur Cayley}
\date{}

\begin{document}
\maketitle

%====================================================================
% PAPER 845 (concluded): ON THE ORTHOMORPHOSIS OF THE CIRCLE INTO
% THE PARABOLA.  pp. 331--336 of printed text.
%====================================================================

\section*{845 (concluded). On the Orthomorphosis of the Circle into the Parabola.}

\medskip

\noindent\textbf{[Printed p.~331]}\hfill

The function $\log\tan(45^{\circ}+\tfrac{1}{2}\arg.)$ is tabulated, Legendre, \textit{Th\'eorie des Fonctions Elliptiques}, t.~II.\ table~iv.\ pp.~256--259, viz.\ writing as above $\phi'$ for the argument, we have hereby the value of $\phi$, $=\log\tan(\tfrac{1}{4}\pi+\tfrac{1}{2}\phi')$, for every value of $\phi'$ from $0^{\circ}$ to $90^{\circ}$ at intervals of $30'$ and to 12 decimals, and with fifth differences. Observe that $\phi'$ is thus given in degrees and minutes as a circular arc, $\phi$ as a number; it is convenient to have $\phi'$ and $\phi$ each as arcs, or each as numbers, the conversion being of course at once made by means of a table of the lengths of circular arcs, and I have calculated the Table which I give at the end of the present paper.

I had previously calculated for the correspondence of $A'O$, $O'B$ and $A'C'B'$ with $AO$, $OB$ and $ACB$, the following values, at irregular intervals suited to the construction of a figure.

\medskip

\begin{center}
\begin{tabular}{|c|c|}
\hline
Circle & Parabola \\
$\angle A'O'P'$ & Ordinate $PM$ \\
\hline
$0^{\circ}$ & $0$ \\
$30$ & $\cdot 3372$ \\
$60$ & $\cdot 6994$ \\
$90$ & $1\cdot 1220$ \\
$120$ & $1\cdot 6768$ \\
$150$ & $2\cdot 5817$ \\
$160$ & $3\cdot 1018$ \\
$170$ & $3\cdot 9869$ \\
$172$ & $4\cdot 2713$ \\
$174$ & $4\cdot 6378$ \\
$176$ & $5\cdot 1542$ \\
$178$ & $6\cdot 0258$ \\
$179$ & $6\cdot 9195$ \\
$180^{\circ}$ & $\infty$ \\
\hline
\end{tabular}
\qquad
\begin{tabular}{|cc|cc|}
\hline
\multicolumn{2}{|c|}{Circle} & \multicolumn{2}{c|}{Parabola} \\
$O'A'$ & $O'B'$ & $OA$ & $OB$ \\
$x'=$ & $x'=$ & $x=$ & $x=$ \\
\hline
$0$ & $0$ & $0$ & $0$ \\
$\cdot 1$ & $-\cdot 1$ & $\cdot 152$ & $-\cdot 173$ \\
$\cdot 2$ & $-\cdot 2$ & $\cdot 287$ & $-\cdot 375$ \\
$\cdot 3$ & $-\cdot 3$ & $\cdot 407$ & $-\cdot 611$ \\
$\cdot 4$ & $-\cdot 4$ & $\cdot 515$ & $-\cdot 943$ \\
$\cdot 5$ & $-\cdot 5$ & $\cdot 615$ & $-1\cdot 257$ \\
$\cdot 6$ & $-\cdot 6$ & $\cdot 704$ & $-1\cdot 724$ \\
$\cdot 7$ & $-\cdot 7$ & $\cdot 787$ & $-2\cdot 368$ \\
$\cdot 8$ & $-\cdot 8$ & $\cdot 853$ & $-3\cdot 386$ \\
$\cdot 9$ & $-\cdot 9$ & $\cdot 935$ & $-5\cdot 368$ \\
$1\cdot 0$ & $-1\cdot 0$ & $1\cdot 000$ & $-\infty$ \\
\hline
\end{tabular}
\end{center}

\medskip

Write in general
\[
\sqrt{(x+iy)} = p+iq,\quad =\frac{a}{\pi}(\psi+i\phi);
\]
viz.\ considering $x$, $y$ as given, find thence $p$, $q$ by the equations $x=p^{2}-q^{2}$, $y=2pq$; and then writing $\tfrac{1}{2}\pi p = \psi$, $\tfrac{1}{2}\pi q = \phi$, we have
\[
\sqrt{(x'+iy')} \;=\; \tan\tfrac{1}{2}\pi(p+iq) \;=\; \tan\bigl(\tfrac{1}{2}\psi+\tfrac{1}{2}i\phi\bigr) \;=\; \frac{P+iQ}{1-iPQ},
\]

\medskip

\noindent\textbf{[Printed p.~332]}\hfill

\noindent where
\[
P = \tan\tfrac{1}{2}\pi p,\quad =\tan\tfrac{1}{2}\psi,
\]
\[
Q = -\tan\tfrac{1}{2}\pi qi,\quad =\tanh\tfrac{1}{2}\pi q,\quad =-\tan\tfrac{1}{2}\phi i,\quad =\tanh\tfrac{1}{2}\phi;
\]
whence, if we introduce $\phi'$ connected with $\phi$ by the equation $\phi=\log\tan\bigl(\tfrac{1}{4}\pi+\tfrac{1}{2}\phi'\bigr)$ as before, we have $Q=\tan\tfrac{1}{2}\phi'$, and the formula is
\[
\sqrt{(x'+iy')} = \frac{P+iQ}{1-iPQ},\quad P=\tan\tfrac{1}{2}\psi,\quad Q=\tan\tfrac{1}{2}\phi',
\]
giving the values of $x'$, $y'$.

It is clear that we have
\[
\sqrt{(x'+iy')} = \frac{P-iQ}{1+iPQ},
\]
and thence
\[
\sqrt{(x'^{2}+y'^{2})} = \frac{P^{2}+Q^{2}}{1+P^{2}Q^{2}}.
\]
Hence, to the circle $x'^{2}+y'^{2}=c^{2}$, corresponds in the parabola the curve
\[
P^{2}+Q^{2} = c\bigl(1+P^{2}Q^{2}\bigr),
\]
where $p+iq=\sqrt{(x+iy)}$, $P=\tan\tfrac{1}{2}\pi p$, $Q=\tanh\tfrac{1}{2}\pi q$. This is a complicated transcendental equation, and I do not see any convenient way of tracing the curve. The set of curves satisfy the differential equation
\[
\frac{P\,dP + Q\,dQ}{P^{2}+Q^{2}} = \frac{PQ(Q\,dP+P\,dQ)}{1+P^{2}Q^{2}},
\]
that is,
\[
P\,dP\,(1-Q^{2}) + Q\,dQ\,(1-P^{2}) = 0,
\]
where $dP$, $dQ$ are given in terms of $dp$, $dq$ by the equations
\[
\frac{dP}{1+P^{2}} = \tfrac{1}{2}\pi\,dp,\qquad \frac{dQ}{1-Q^{2}} = \tfrac{1}{2}\pi\,dq.
\]

We have $y=2pq$, and thence at the highest points, or summits of the several curves, $q\,dp + p\,dq = 0$. Combining these equations, we have
\[
dP : dQ = p(1+P^{2}) : -q(1-Q^{2}),
\]
and thence
\[
p\,P\,(1+P^{2})(1-Q^{2}) - q\,Q\,(1-P^{2})(1-Q^{2}) = 0,
\]
as an equation to the locus of the summits in question. If $p$ and $q$ are small, then putting for a moment $\tfrac{1}{2}\pi = m$ for shortness, we have
\[
P = mp + \tfrac{1}{3}m^{3}p^{3},\quad Q = mq - \tfrac{1}{3}m^{3}q^{3},
\]
and the equation becomes
\[
p^{2}(1+m^{2}p^{2})(1+m^{2}q^{2}) - q^{2}(1-m^{2}p^{2})(1-m^{2}q^{2}) = 0,
\]

\medskip

\noindent\textbf{[Printed p.~333]}\hfill

\noindent that is,
\[
p^{2} - q^{2} + \tfrac{1}{3}m^{2}(p^{4} + 6p^{2}q^{2} + q^{4}) = 0;
\]
writing this in the form
\[
p^{2} - q^{2} + \tfrac{m^{2}}{3}\bigl\{(p^{2}-q^{2})^{2} + 8p^{2}q^{2}\bigr\} = 0,
\]
we find $1=p^{2}+q^{2}$, $\tfrac{4m^{2}}{3}=-\tfrac{x}{x^{2}-y^{2}}$, or say $x^{2}+\tfrac{3}{2m^{2}}x - y^{2}=0$, as the locus of the summits in the neighbourhood of the focus $O$, viz.\ the summits lie all of them, as might have been expected, on the left-hand (or negative side) of the focus.

I have constructed for the parabola, to the scale $1=1\tfrac{1}{4}$ inch, as accurately as the data enable, the figure corresponding to the concentric circles and the radii of the circle.

Resuming the equation $\sqrt{(x+iy)}=p+iq$, that is, $x=p^{2}-q^{2}$ and $y=2pq$, we have the curves $(p=\textnormal{const.})$, $y^{2}=4p^{2}(p^{2}-x)$, and $(q=\textnormal{const.})$, $y^{2}=4q^{2}(q^{2}+x)$, which are two systems of confocal parabolas, cutting each other at right angles; the curves of the former set all of them interior to the given parabola, those of the latter set of course cutting it at right angles.

Corresponding hereto in the circle, writing for a moment
\[
\sqrt{(x'+iy')} = p'+iq',
\]
we have
\[
p'+iq' = \frac{P+iQ}{1-iPQ},
\]
whence
\[
p' = P(1-q'Q),\quad q' = Q(1+p'P);
\]
whence eliminating $P$ and $Q$ successively, we find
\[
p'^{2} + q'^{2} - p'\left(P - \frac{1}{P}\right) - 1 = 0,
\]
\[
p'^{2} + q'^{2} - q'\left(Q + \frac{1}{Q}\right) + 1 = 0,
\]
say, for shortness, these are $p'^{2} + q'^{2} - 2mp' + 1 = 0$, and $p'^{2} + q'^{2} - 2nq' - 1 = 0$. Introducing the polar coordinates $r'$, $\theta'$, we have
\[
x' + iy' = r'(\cos\theta' + i\sin\theta'),
\]
and thence
\[
p' = \sqrt{(r')}\cos\tfrac{1}{2}\theta',\quad q' = \sqrt{(r')}\sin\tfrac{1}{2}\theta';
\]
the equations thus become
\[
r' - 1 - 2m\sqrt{(r')}\cos\tfrac{1}{2}\theta' = 0,
\]
\[
r' + 1 - 2n\sqrt{(r')}\sin\tfrac{1}{2}\theta' = 0,
\]

\medskip

\noindent\textbf{[Printed p.~334]}\hfill

\noindent which belong to two lima\c{c}ons having each of them $B'$ for a node and $O'$ for a focus, and which of course cut each other at right angles; see fig.~4, which is a mere diagram.

\medskip

\textit{[Figure 4: a diagram showing two orthogonal lima\c{c}ons inscribed in a circle, with labelled points $B'$, $O'$, $A'$ on the horizontal diameter.]}

\medskip

In fact, omitting for convenience the accents, but recollecting always that the curves belong to the figure of the circle, the first equation gives
\[
(r-1)^{2} = 4m^{2}r\cos^{2}\tfrac{1}{2}\theta,\quad = 2m^{2}r(1+\cos\theta) = 2m^{2}(r+x),
\]
that is,
\[
(r^{2}+1-2m^{2}x)^{2} = 4(m^{2}+1)^{2} r^{2}.
\]

Transforming to the point $B'$ as origin, we must write $x+1$ for $x$, and then the equation thus becomes
\[
\bigl\{x^{2}+y^{2}-2(m^{2}+1)x + 2(m^{2}+1)\bigr\}^{2} = 4(m^{2}+1)^{2}(x^{2}+y^{2}+2x+1),
\]
that is,
\[
\bigl\{x^{2}+y^{2}-2(m^{2}+1)x\bigr\}^{2} = 4m^{2}(m^{2}+1)(x^{2}+y^{2}),
\]
or say
\[
(x^{2}+y^{2})^{2} - 4(m^{2}+1)(x^{2}+y^{2})x + 4(m^{2}+1)(x^{2}-m^{2}y^{2}) = 0,
\]
showing that the point $B'$ is a crunode, the tangents there being $x=\pm my$.

Writing in the equation $y=0$, we have $x=0$, the crunode, and
\[
x = 2\sqrt{(m^{2}+1)}\bigl\{\sqrt{(m^{2}+1)}\pm m\bigr\};
\]
the product of these two values is $=4(m^{2}+1)$, that is, $=\left(P+\tfrac{1}{P}\right)^{2}$, viz.\ one value is greater, the other less, than $2$. We see also that
\[
2\sqrt{(m^{2}+1)}\bigl\{\sqrt{(m^{2}+1)}-m\bigr\},
\]

\medskip

\noindent\textbf{[Printed p.~335]}\hfill

\noindent the smaller root, is greater than $1$. The curve corresponding to a parabola $y^{2}=4p^{2}(p^{2}-x)$ is thus a crunodal lima\c{c}on, the crunode at $B'$, and the loop lying wholly within the circle. Moreover, the loop includes within itself the centre $O'$ of the circle.

The other curve is in like manner
\[
\bigl\{r^{2}+1+2n^{2}x\bigr\}^{2} = 4(n^{2}+1)r^{2},
\]
viz.\ transforming to the origin $B'$, and therefore putting $x-1$ for $x$, and
\[
r^{2} = (x-1)^{2}+y^{2},
\]
the equation is
\[
\bigl\{x^{2}+y^{2}-2(n^{2}+1)x + 2(n^{2}+1)\bigr\}^{2} = 4(n^{2}+1)^{2}(x^{2}+y^{2}-2x+1),
\]
that is,
\[
\bigl\{x^{2}+y^{2}-2(n^{2}+1)x\bigr\}^{2} = 4n^{2}(n^{2}+1)(x^{2}+y^{2}),
\]
or say
\[
(x^{2}+y^{2})^{2} - 4(n^{2}+1)(x^{2}+y^{2})x - 4(n^{2}+1)(x^{2}+n^{2}y^{2}) = 0,
\]
viz.\ the point $B'$ is an acnode with the imaginary tangents $x=\pm iny$.

Writing in the equation $y=0$, we have $x=0$ the acnode, and
\[
x = -2\sqrt{(n^{2}-1)}\bigl\{\sqrt{(n^{2}-1)}\pm n\bigr\},
\]
where
\[
\sqrt{(n^{2}-1)} = \tfrac{1}{2}\sqrt{\left(Q-\tfrac{1}{Q}\right)^{2}}
\]
is real and may be taken to be positive; there is thus one root
\[
-\sqrt{(n^{2}-1)}\bigl\{\sqrt{(n^{2}-1)}+n\bigr\},
\]
which is negative; the other root, say
\[
\sqrt{(n^{2}-1)}\bigl\{n - \sqrt{(n^{2}-1)}\bigr\},
\]
is positive and less than $1$; the curve is thus an acnodal lima\c{c}on, having $B'$ for the acnode and cutting the axis outside the circle on the negative side of $B'$, and inside the circle between $B'$ and $O'$.

Taking $B'$ as origin, the equation of the circle is $x^{2}+y^{2}-2x=0$, and we hence find for the intersection of the lima\c{c}on with the circle $n^{2}x = 2(n^{2}-1)$, that is,
\[
x = 2\left(1-\frac{1}{n^{2}}\right),\quad y = \pm\sqrt{\left(\frac{4}{n^{2}}-\frac{1}{n^{4}}\right)};
\]
whence also
\[
r^{2} = 4\left(1-\frac{1}{n^{2}}\right)^{2} + \frac{4}{n^{2}} - \frac{1}{n^{4}},
\]
values which belong to a real intersection; for $q=0$, we have $Q=0$, $n=\infty$, and therefore $(x,y)=(2,0)$, viz.\ the intersection is at $A'$; for $q=\infty$, we have $Q=1$, $n=1$, and therefore $(x,y)=(0,0)$, viz.\ the intersection is at $B'$, results which are obviously right. Observe that for the other lima\c{c}on we have
\[
x = 2\left(1+\frac{1}{m^{2}}\right),\quad y = \pm\sqrt{\left(-\frac{4}{m^{2}} - \frac{1}{m^{4}}\right)},
\]
viz.\ there is no real intersection with the circle.

\medskip

\noindent\textbf{[Printed p.~336]}\hfill

\begin{center}
\textit{The Table above referred to.}\\[2pt]
$\phi = \log\tan\bigl(45^{\circ}+\tfrac{1}{2}\phi'\bigr)$.
\end{center}

\begin{center}
\small
\begin{tabular}{|r|r|r||r|r|r|}
\hline
$\phi'$ & $\phi$ & $\phi$ (arc) & $\phi'$ & $\phi$ & $\phi$ (arc) \\
\hline
$0^{\circ}$ & $\cdot 0$ & $\cdot 0$ & $46^{\circ}$ & $\cdot 8028515$ & $\cdot 9062755$ \\
$1$ & $\cdot 0174533$ & $\cdot 0174542$ & $47$ & $\cdot 8203047$ & $\cdot 9316316$ \\
$2$ & $\cdot 0349066$ & $\cdot 0349137$ & $48$ & $\cdot 8377580$ & $\cdot 9574669$ \\
$3$ & $\cdot 0523599$ & $\cdot 0523838$ & $49$ & $\cdot 8552113$ & $\cdot 9838079$ \\
$4$ & $\cdot 0698132$ & $\cdot 0698699$ & $50$ & $\cdot 8726646$ & $1\cdot 0106832$ \\
$5$ & $\cdot 0872665$ & $\cdot 0873774$ & $51$ & $\cdot 8901179$ & $1\cdot 0381235$ \\
$6$ & $\cdot 1047198$ & $\cdot 1049117$ & $52$ & $\cdot 9075712$ & $1\cdot 0661617$ \\
$7$ & $\cdot 1221730$ & $\cdot 1224781$ & $53$ & $\cdot 9250245$ & $1\cdot 0948335$ \\
$8$ & $\cdot 1396263$ & $\cdot 1400822$ & $54$ & $\cdot 9424778$ & $1\cdot 1241772$ \\
$9$ & $\cdot 1570796$ & $\cdot 1577296$ & $55$ & $\cdot 9509311$ & $1\cdot 1542316$ \\
$10$ & $\cdot 1745329$ & $\cdot 1754258$ & $56$ & $\cdot 9773844$ & $1\cdot 1850507$ \\
$11$ & $\cdot 1919862$ & $\cdot 1931766$ & $57$ & $\cdot 9948377$ & $1\cdot 2166748$ \\
$12$ & $\cdot 2094395$ & $\cdot 2109867$ & $58$ & $1\cdot 0122910$ & $1\cdot 2491606$ \\
$13$ & $\cdot 2268928$ & $\cdot 2288650$ & $59$ & $1\cdot 0297443$ & $1\cdot 2825668$ \\
$14$ & $\cdot 2443461$ & $\cdot 2468145$ & $60$ & $1\cdot 0471976$ & $1\cdot 3169579$ \\
$15$ & $\cdot 2617994$ & $\cdot 2648422$ & $61$ & $1\cdot 0646508$ & $1\cdot 3524048$ \\
$16$ & $\cdot 2792527$ & $\cdot 2829545$ & $62$ & $1\cdot 0821041$ & $1\cdot 3889860$ \\
$17$ & $\cdot 2967060$ & $\cdot 3011577$ & $63$ & $1\cdot 0995574$ & $1\cdot 4267882$ \\
$18$ & $\cdot 3141593$ & $\cdot 3194583$ & $64$ & $1\cdot 1170107$ & $1\cdot 4659083$ \\
$19$ & $\cdot 3316126$ & $\cdot 3378629$ & $65$ & $1\cdot 1344640$ & $1\cdot 5064542$ \\
$20$ & $\cdot 3490659$ & $\cdot 3563785$ & $66$ & $1\cdot 1519173$ & $1\cdot 5485472$ \\
$21$ & $\cdot 3665191$ & $\cdot 3750121$ & $67$ & $1\cdot 1693706$ & $1\cdot 5923237$ \\
$22$ & $\cdot 3839724$ & $\cdot 3937710$ & $68$ & $1\cdot 1868239$ & $1\cdot 6379387$ \\
$23$ & $\cdot 4014257$ & $\cdot 4126626$ & $69$ & $1\cdot 2042772$ & $1\cdot 6855685$ \\
$24$ & $\cdot 4188790$ & $\cdot 4316947$ & $70$ & $1\cdot 2217305$ & $1\cdot 7354152$ \\
$25$ & $\cdot 4363323$ & $\cdot 4508753$ & $71$ & $1\cdot 2391838$ & $1\cdot 7877120$ \\
$26$ & $\cdot 4537856$ & $\cdot 4702127$ & $72$ & $1\cdot 2566371$ & $1\cdot 8427300$ \\
$27$ & $\cdot 4712389$ & $\cdot 4897154$ & $73$ & $1\cdot 2740904$ & $1\cdot 9007867$ \\
$28$ & $\cdot 4886922$ & $\cdot 5093923$ & $74$ & $1\cdot 2915436$ & $1\cdot 9622572$ \\
$29$ & $\cdot 5061455$ & $\cdot 5292527$ & $75$ & $1\cdot 3089969$ & $2\cdot 0275894$ \\
$30$ & $\cdot 5235988$ & $\cdot 5493061$ & $76$ & $1\cdot 3264502$ & $2\cdot 0973240$ \\
$31$ & $\cdot 5410521$ & $\cdot 5695627$ & $77$ & $1\cdot 3439035$ & $2\cdot 1721218$ \\
$32$ & $\cdot 5585054$ & $\cdot 5900329$ & $78$ & $1\cdot 3613568$ & $2\cdot 2528027$ \\
$33$ & $\cdot 5759587$ & $\cdot 6107275$ & $79$ & $1\cdot 3788101$ & $2\cdot 3404007$ \\
$34$ & $\cdot 5934119$ & $\cdot 6316581$ & $80$ & $1\cdot 3962634$ & $2\cdot 4362460$ \\
$35$ & $\cdot 6108652$ & $\cdot 6528366$ & $81$ & $1\cdot 4137167$ & $2\cdot 5420904$ \\
$36$ & $\cdot 6283185$ & $\cdot 6742755$ & $82$ & $1\cdot 4311700$ & $2\cdot 6603061$ \\
$37$ & $\cdot 6457718$ & $\cdot 6959880$ & $83$ & $1\cdot 4486233$ & $2\cdot 7942190$ \\
$38$ & $\cdot 6632251$ & $\cdot 7179880$ & $84$ & $1\cdot 4660766$ & $2\cdot 9487002$ \\
$39$ & $\cdot 6806784$ & $\cdot 7402901$ & $85$ & $1\cdot 4835299$ & $3\cdot 1313013$ \\
$40$ & $\cdot 6981317$ & $\cdot 7629095$ & $86$ & $1\cdot 5009832$ & $3\cdot 3546735$ \\
$41$ & $\cdot 7155850$ & $\cdot 7858630$ & $87$ & $1\cdot 5184364$ & $3\cdot 6425334$ \\
$42$ & $\cdot 7330383$ & $\cdot 8091672$ & $88$ & $1\cdot 5358897$ & $4\cdot 0481254$ \\
$43$ & $\cdot 7504916$ & $\cdot 8328406$ & $89$ & $1\cdot 5533430$ & $4\cdot 7413488$ \\
$44$ & $\cdot 7679449$ & $\cdot 8569026$ & $90$ & $1\cdot 5707963$ & $\infty$ \\
$45$ & $\cdot 7853982$ & $\cdot 8813736$ & & & \\
\hline
\end{tabular}
\end{center}

\newpage

%====================================================================
% PAPER 846.  A VERIFICATION IN REGARD TO THE LINEAR TRANSFORMATION
% OF THE THETA-FUNCTIONS.  pp. 337--343.
%====================================================================

\section*{846. A Verification in Regard to the Linear Transformation of the Theta-Functions.}

\noindent\textit{[From the Quarterly Journal of Pure and Applied Mathematics, vol.~xxi.\ (1886), pp.~77--84.]}

\medskip

\noindent\textbf{[Printed p.~337]}\hfill

The notation is that of Smith's* ``Memoir on the Theta and Omega Functions,'' differing from Hermite's only in a factor, $\vartheta_{m,n}(\mathrm{Smith}) = i^{mn}\vartheta_{m,n}(\mathrm{Hermite})$; and it is in consequence of this that the factor $i^{\frac{1}{2}\nu - mn}$ occurs in the expression presently given for $\mathcal{C}$. Hermite's Memoir ``Sur quelques formules relatives \`a la transformation des fonctions elliptiques'' appeared in Liouville, t.~iii.\ (1858), pp.~27--36.

\medskip

\footnotetext{[* H.~J.~S.~Smith, \textit{Collected Mathematical Papers}, vol.~II., pp.~415--621.]}

Writing
\[
\omega = \frac{c + d\Omega}{a + b\Omega},
\]
where $ad - bc = 1$, the formula of transformation is
\[
\vartheta_{m,n}\!\left\{(a+b\Omega)\frac{\pi x}{h},\ \Omega\right\} \;=\; \mathcal{C}\,\exp\!\left\{-i\pi b(a+b\Omega)\frac{x^{2}}{h^{2}}\right\}\vartheta_{m',n'}\!\left(\frac{\pi x}{h},\ \omega\right),
\]
where
\[
m' = a\mu + b\nu + ab,
\]
\[
n' = c\mu + d\nu + cd.
\]

We have, according as $b$ is positive or negative,
\[
C = \frac{\delta H}{\sqrt{\{-i(a+b\Omega)\}}},\quad\text{or}\quad C = \frac{\delta H}{\sqrt{\{i(a+b\Omega)\}}},
\]
where for each case the square root is to be taken in such wise that the real part is positive; $\delta$, $H$ are eighth roots of unity; viz.\ in each case
\[
\delta = \exp\bigl\{-\tfrac{1}{4}i\pi\bigl(ac\mu^{2} + 2bc\mu\nu + bd\nu^{2} - 2abc\mu - 2abd\nu + ab^{2}c\bigr)\bigr\}\, i^{\tfrac{1}{2}\nu - mn},
\]

\medskip

\noindent\textbf{[Printed p.~338]}\hfill

\noindent and, according as $b$ is positive or negative,
\[
H = \left(\frac{a}{b}\right) i^{-\frac{1}{2}a},\quad a\ \text{odd}, \qquad H = \left(\frac{-b}{a}\right) \beta^{n},\quad a\ \text{odd},
\]
\[
H = \left(\frac{a}{b}\right) i^{\frac{1}{2}a + \frac{1}{2}(a - b)},\quad b\ \text{odd}, \qquad H = \left(\frac{-a}{b}\right) \beta^{n - \frac{1}{2}a + \frac{1}{2}(a + b)},\quad b\ \text{odd},
\]
where $\left(\dfrac{a}{b}\right)$, \&c., are the Legendrian symbols as generalised by Jacobi; if $a$ and $b$ are each of them odd, the formulae for the cases $a$ odd and $b$ odd respectively are equivalent to each other, and either may be used indifferently.

To fix the ideas, I assume throughout that $b$ is positive and odd; the proper formulae thus are
\[
C = \frac{\delta H}{\sqrt{\{-i(a+b\Omega)\}}},
\]
$\delta$ as above, and
\[
H = \left(\frac{a}{b}\right) i^{\frac{1}{2}a + \frac{1}{2}(a - b)}.
\]

I will also, for greater convenience, assume that $c$ is odd; $ad$ is consequently even, viz.\ the numbers $a$ and $d$ are each of them even, or else one is odd and the other even.

The equation $\omega = \dfrac{c + d\Omega}{a + b\Omega}$ gives $\Omega = \dfrac{c - a\omega}{-d + b\omega}$, and thence
\[
a + b\Omega = \frac{-1}{-d + b\omega}.
\]

Comparing the expressions for $\omega$, $\Omega$, it appears that we may interchange $\omega$ and $\Omega$, writing also $-d$, $b$, $-c$, for $a$, $b$, $c$, $d$; and changing the other letters, we may for
\[
a,\ b,\ c,\ d,\ \Omega,\ \mu,\ \nu,\ \omega,\ m',\ n',\ \text{write}
\]
\[
-d,\ b,\ -c,\ a,\ \omega,\ m',\ n',\ \Omega,\ m,\ n,
\]
where
\[
m' = -dm + bn - bd,
\]
\[
n' = -cm + an - ac.
\]

The formula thus becomes
\[
\vartheta_{m,n}\!\left\{-\frac{\pi x}{h},\ \omega\right\} = C'\,\exp\!\left\{-i\pi(-d+b\omega)\frac{x^{2}}{h^{2}}\right\}\vartheta_{m',n'}\!\left(\frac{\pi x}{h},\ \Omega\right),
\]
or, for $x$ writing $(a+b\Omega)x$, this is
\[
\vartheta_{m,n}\!\left\{-\frac{\pi x}{h}(a+b\Omega),\ \omega\right\} = C'\,\exp\!\left\{-i\pi(a+b\Omega)\frac{x^{2}}{h^{2}}\right\}\vartheta_{m',n'}\!\left((a+b\Omega)\frac{\pi x}{h},\ \Omega\right),
\]

\medskip

\noindent\textbf{[Printed p.~339]}\hfill

\noindent or, observing that the left-hand side is $=(-)^{mn}\vartheta_{m,n}\!\left\{(a+b\Omega)\dfrac{\pi x}{h},\ \omega\right\}$, we have
\[
\vartheta_{m,n}\!\left\{(a+b\Omega)\frac{\pi x}{h},\ \omega\right\} = \frac{(-)^{mn}}{C'}\exp\!\left\{-i\pi(a+b\Omega)\frac{x^{2}}{h^{2}}\right\}\vartheta_{m',n'}\!\left(\frac{\pi x}{h},\ \Omega\right),
\]
an equation which is of the same form as the original equation of transformation, and is to be identified with it.

We have
\begin{align*}
m' &= -dm + bn - bd = -d(a\mu + b\nu + ab) + b(c\mu + d\nu + cd) - bd \\
   &= \mu + bd(-a + c - 1), \\
n' &= -cm + an - ac = -c(a\mu + b\nu + ab) + a(c\mu + d\nu + cd) - ac \\
   &= \nu + ac(-b + d - 1);
\end{align*}
values which may be written
\[
m' = 2P - \mu,\quad \text{where } P = \tfrac{1}{2}bd(-a + c - 1),
\]
\[
n' = 2Q - \nu,\quad \text{where } Q = \tfrac{1}{2}ac(-b + d - 1),
\]
$P$ and $Q$ being integers. In fact, $P$ is an integer if $b$ or $d$ is even, and if they are each of them odd, then from the equation $ad - bc = 1$, $a$ and $c$ will be one of them odd, the other even, and $-a + c - 1$ will be even; and similarly for $Q$.

The left-hand side is
\[
\vartheta_{m',n'}\!\left\{(a+b\Omega)\frac{\pi x}{h},\ \Omega\right\},
\]
which is
\[
= (-)^{Pv}\vartheta_{-\mu,-\nu}\!\left\{(a+b\Omega)\frac{\pi x}{h},\ \Omega\right\} = (-)^{Pv}\vartheta_{\mu,\nu}\!\left\{(a+b\Omega)\frac{\pi x}{h},\ \Omega\right\},
\]
and the equation finally is
\[
\vartheta_{\mu,\nu}\!\left\{(a+b\Omega)\frac{\pi x}{h},\ \Omega\right\} = \frac{1}{C'}(-)^{mn + Pv}\exp\!\left\{-i\pi(a+b\Omega)\frac{x^{2}}{h^{2}}\right\}\vartheta_{m,n}\!\left(\frac{\pi x}{h},\ \omega\right).
\]

Comparing this with the original equation
\[
\vartheta_{\mu,\nu}\!\left\{(a+b\Omega)\frac{\pi x}{h},\ \Omega\right\} = C\,\exp\!\left\{-i\pi(a+b\Omega)\frac{x^{2}}{h^{2}}\right\}\vartheta_{m,n}\!\left(\frac{\pi x}{h},\ \omega\right),
\]
the two equations will be identical if only
\[
CC' = (-)^{-mn - Pv}.
\]

We have
\[
C = \frac{\delta H}{\sqrt{\{-i(a+b\Omega)\}}},\qquad C' = \frac{\delta' H'}{\sqrt{\{-i(-d+b\omega)\}}},
\]
where the square roots are taken in such wise that the real part is positive; hence
\[
CC' = \frac{\delta\delta' H H'}{\sqrt{(+1)}},
\]

\medskip

\noindent\textbf{[Printed p.~340]}\hfill

\noindent where $\sqrt{(+1)}$ means
\[
\sqrt{\{-i(a+b\Omega)\}}\cdot\sqrt{\{-i(-d+b\omega)\}},
\]
the last-mentioned two square roots being as just explained; and we have moreover
\[
\delta = \exp\bigl\{-\tfrac{1}{4}i\pi\bigl(ac\mu^{2} + 2bc\mu\nu + bd\nu^{2} - 2abc\mu - 2abd\nu + ab^{2}c\bigr)\bigr\}\, i^{\tfrac{1}{2}\nu - mn},
\]
\[
\delta' = \exp\bigl\{-\tfrac{1}{4}i\pi\bigl(-dcm^{2} + 2bcmn - abn^{2} - 2dbcm + 2abdn - db^{2}c\bigr)\bigr\}\, i^{\tfrac{1}{2}n' - m'n'},
\]
viz.\ the value of $\delta'$ is obtained from that of $\delta$ by the change $a$, $b$, $c$, $d$, $\mu$, $\nu$, $m$, $n$ into $-d$, $b$, $-c$, $a$, $m$, $n$, $m'$, $n'$.

Representing these by
\[
\delta = \exp\bigl\{-\tfrac{1}{4}(i\pi)\,\Delta\bigr\}\, i^{\tfrac{1}{2}\nu - mn},\qquad \delta' = \exp\bigl\{-\tfrac{1}{4}(i\pi)\,\Delta'\bigr\}\, i^{\tfrac{1}{2}n' - m'n'},
\]
we have
\[
\delta\delta' = \exp\bigl\{-\tfrac{1}{4}i\pi(\Delta + \Delta')\bigr\}\, i^{\tfrac{1}{2}(\nu + n') - mn - m'n'}.
\]

But
\[
m'n' - mn = (2P - \mu)(2Q - \nu) - \mu\nu = 4PQ - 2P\nu - 2Q\mu + \mu\nu - \mu\nu,
\]
that is,
\[
mn - m'n' = -4PQ + 2P\nu + 2Q\mu,
\]
or, omitting the term divisible by $4$,
\[
i^{2P\nu + 2Q\mu} = (-)^{P\nu + Q\mu},
\]
\[
i^{\tfrac{1}{2}(\nu + n') - mn - m'n'} = i^{2P\nu + 2Q\mu} \cdot (-)^{P\nu + Q\mu}.
\]

To calculate $\Delta + \Delta'$, we have
\[
dm - bn = \mu + bd(-a + c),
\]
\[
-cm + an = \nu + ac(-b + d),
\]
and thence
\begin{align*}
-cdm^{2} + (ad + bc)mn - abn^{2} &= \mu\nu + \mu ac(-b + d) + \nu bd(a - c) + abcd(d - b)(a - c) \\
&\quad + ac\mu^{2} - bd\nu^{2} - (ad + bc)\mu\nu - \mu ac(b - d) - \nu bd(a - c) \\
(-ad + bc)\,mn &= - abcd,
\end{align*}
consequently
\[
-cdm^{2} + 2bcmn - abn^{2} = -ac\mu^{2} - 2bc\mu\nu - bd\nu^{2} - 2abc\mu - 2bcd\nu + abcd(2bc - ab - cd);
\]
also
\begin{align*}
-2bcdm &= -2abcd\mu - 2b^{2}cd\nu - 2ab^{2}cd, \\
2abdn &= -2abcd\mu + 2abd^{2}\nu + 2abcd^{2}, \\
-db^{2}c &= -db^{2}c;
\end{align*}
whence, adding, we obtain
\begin{align*}
\Delta' = -ac\mu^{2} - 2bc\mu\nu - bd\nu^{2} - 2abc\mu &+ (-2bcd - 2b^{2}cd - 2abd^{2})\nu \\
 &+ abcd(2bc - ab - cd - 2b + 2d) - db^{2}c,
\end{align*}
and, adding to this the expression of $\Delta$, we find
\[
\Delta + \Delta' = (2abd - 2bcd - 2b^{2}cd + 2abd^{2})\nu - ab^{2}c - db^{2}c + abcd(2bc - ab - cd - 2b + 2d).
\]

\medskip

\noindent\textbf{[Printed p.~341]}\hfill

\noindent The coefficient of $\nu$ is $= 2bd(a - c - bc + ad) = 2bd(a + c - 1)$, which is $= -4P\nu$. Hence writing
\[
\Theta = \tfrac{1}{2}abcd(2bc - ab - cd - 2b + 2d) + \tfrac{1}{2}b^{2}c(a - d),
\]
where observe that $4\Theta$, but not in every case $\Theta$ itself, is an integer, we have
\[
\Delta + \Delta' = -4P\nu + 4\Theta,
\]
and consequently
\[
\delta\delta' = (-)^{-P\nu + \Theta + Pv + Q\mu} \cdot i^{2P\nu + 2Q\mu} = (-)^{\Theta + P\nu + Q\mu} \cdot i^{2P\nu + 2Q\mu}.
\]
Observe that $(-)$ denotes, and it might properly have been written, $\exp i\pi\Theta$. The foregoing equation
\[
\frac{\delta\delta'\,HH'}{\sqrt{(+1)}} = (-)^{-mn - Pv}
\]
becomes thus
\[
\frac{HH'}{\sqrt{(+1)}} = (-)^{-mn - Pv} \cdot \frac{1}{\delta\delta'},
\]
that is,
\[
\frac{HH'}{\sqrt{(+1)}} = (-)^{-mn - Pv - \Theta - P\nu - Q\mu} \cdot i^{-2P\nu - 2Q\mu},
\]
where the even term $2Q\mu$ may be omitted. We have moreover
\begin{align*}
mn &= ac\mu^{2} + (ad + bc)\mu\nu + bd\nu^{2} + ac(b + d)\mu + bd(a - c)\nu + abcd, \\
\mu\nu &= (-ad + bc)\mu\nu,
\end{align*}
and thence
\[
mn + \mu\nu = ac(\mu^{2} + \mu) + 2bc\mu\nu + bd(\nu^{2} + \nu) + ac(b + d - 1)\mu + bd(a - c - 1)\nu + abcd,
\]
where each term is even; hence $mn - \mu\nu$ is even, and we have simply
\[
\frac{HH'}{\sqrt{(+1)}} = (-)^{-\mu\nu - B},
\]
where
\[
B = \tfrac{1}{2}abcd(ab + cd - 2bc + 2b - 2d) - \tfrac{1}{2}b^{2}c(a - d).
\]

If I write $M = \tfrac{1}{2}b(a - d)$, then
\[
\Theta - M = \tfrac{1}{2}abcd(ab + cd - 2bc + 2b - 2d) - \tfrac{1}{2}(bc + 1)b(a - d),
\]
where the second term is $= \tfrac{1}{2}abd(a - d)$, and we have therefore
\[
\Theta - M = \tfrac{1}{2}abd(abc + c^{2}d - a + 2bc^{2} - 2bc - 2cd - a + d).
\]

I assume, as above, that $b$ and $c$ are each of them odd; therefore $ad$ is even. I suppose, first, that $ad$ divides by $4$, then $\tfrac{1}{2}abd$ is an integer, and in the expression of $\Theta - M$, omitting even numbers, we have
\[
\tfrac{1}{2}abd(abc + c^{2}d - a + d),
\]

\medskip

\noindent\textbf{[Printed p.~342]}\hfill

\noindent which, putting therein $bc = ad - 1$, becomes
\begin{align*}
&= \tfrac{1}{2}abd(a^{2}d + c^{2}d - 2a + d), \\
&= \tfrac{1}{2}abd\bigl\{a^{2}d + (c^{2} - c)d + (c + 1)d - 2a\bigr\},
\end{align*}
where inside the $\{\ \}$ each term is even; hence $\Theta - M$ is even.

Next, if $ad$ is even but not divisible by $4$, then $bc = ad - 1$, which is $\equiv 1\ (\modd 4)$; thus $b$ and $c$ are
\[
= 4\sigma + 1\quad\text{and}\quad 4\tau + 1,
\]
or else
\[
= 4\sigma - 1\quad\text{and}\quad 4\tau - 1,
\]
and, moreover, $bc = 4\theta + 1$, and $c^{2} = 4\phi + 1$; hence
\[
\Theta - M = \tfrac{1}{2}abd\bigl\{a(4\theta + 1) + (4\phi + 1)d - 2a + 2(4\theta + 1)d - 2a + d\bigr\},
\]
or, omitting even numbers, that is, inside the $\{\ \}$ numbers which contain the factor $4$, this is
\[
= \tfrac{1}{2}abd(a + d - 2b + 2 - a + d),
\]
\[
= \tfrac{1}{2}abd(d - b + 1 - cd),
\]
or, since each term within the $\{\ \}$ is even, we have in this case also $\Theta - M$ even.

And this being so, the foregoing equation for $HH'$ becomes
\[
\frac{HH'}{\sqrt{(+1)}} = (-)^{-M},\quad\text{or say } = i^{2M}.
\]

The values of $H$, $H'$, refer each of them to a positive odd value of $b$, and they thus are
\[
H = \left(\frac{a}{b}\right) i^{\frac{1}{2}a + \frac{1}{2}(a - b)},\quad H' = \left(\frac{-d}{b}\right) i^{\frac{1}{2}d + \frac{1}{2}(-d - b)};
\]
hence
\[
HH' = \left(\frac{a}{b}\right)\left(\frac{-d}{b}\right) i^{\frac{1}{2}a + \frac{1}{2}d + (a - b)} = \left(\frac{a}{b}\right)\left(\frac{-d}{b}\right) i^{\frac{1}{2}(a + d) - 2 + 2(a - b)},
\]
or, since
\[
\left(\frac{-d}{b}\right) = (-)^{\frac{1}{2}(b - 1)}\left(\frac{d}{b}\right) = i^{2 \cdot \frac{1}{2}(b - 1)}\left(\frac{d}{b}\right),
\]
and $2(b - 1)$ divides by $4$, this is
\[
HH' = \left(\frac{a}{b}\right)\left(\frac{d}{b}\right) i^{\frac{1}{2}(a + d)}.
\]
Also $\left(\dfrac{a}{b}\right)\left(\dfrac{d}{b}\right) = \left(\dfrac{ad}{b}\right)$, but from the equation $ad - bc = 1$, or $\dfrac{ad}{b} = c + \dfrac{1}{b}$, we have $\left(\dfrac{ad}{b}\right) = \left(\dfrac{1}{b}\right) = 1$; whence
\[
HH' = i^{\frac{1}{2}(a + d)}.
\]

\medskip

\noindent\textbf{[Printed p.~343]}\hfill

We have $\omega = x + iy$, where $y$ is positive; hence
\[
-d + b\omega = -d + bx + iby = \alpha + iby,\quad \text{if } \alpha = -d + bx;
\]
hence
\begin{align*}
a + b\Omega &= \frac{-1}{\alpha + iby} = \frac{-\alpha + iby}{\alpha^{2} + \beta^{2}y^{2}},\quad \beta = b, \\
-i(a + b\Omega) &= \frac{by - i\alpha}{\alpha^{2} + \beta^{2}y^{2}} = R(\cos\theta + i\sin\theta), \\
-i(-d + b\Omega) &= by - i\alpha = R(\cos\theta + i\sin\theta), \\
-i(a + b\Omega) &= \frac{by - i\alpha}{\alpha^{2} + \beta^{2}y^{2}} = \frac{R(\cos\theta - i\sin\theta)}{\alpha^{2} + \beta^{2}y^{2}},
\end{align*}
where $R$ is positive; and $\cos\theta$ is positive since $y$ is positive, and thus $\theta$ lies between $-\tfrac{1}{2}\pi$ and $\tfrac{1}{2}\pi$. Hence
\[
\sqrt{\{-i(-d + b\Omega)\}} = \sqrt{R}(\cos\tfrac{1}{2}\theta + i\sin\tfrac{1}{2}\theta),
\]
$\cos\tfrac{1}{2}\theta$ being positive,
\[
\sqrt{\{-i(a + b\Omega)\}} = \frac{\sqrt{R}(\cos\tfrac{1}{2}\theta - i\sin\tfrac{1}{2}\theta)}{\sqrt{(\alpha^{2} + \beta^{2}y^{2})}},
\]
and thus
\[
\sqrt{\{-i(a + b\Omega)\}}\cdot\sqrt{\{-i(-d + b\Omega)\}} = \frac{R}{\sqrt{(\alpha^{2} + \beta^{2}y^{2})}} = +1,
\]
that is, $\sqrt{(+1)} = +1$, and we thus have, as we should do,
\[
\frac{HH'}{\sqrt{(+1)}} = i^{\frac{1}{2}(a + d)},
\]
the equation which was to be verified.

\newpage

%====================================================================
% PAPER 847. ON THE THEORY OF SEMINVARIANTS.  pp. 344-- (start)
%====================================================================

\section*{847. On the Theory of Seminvariants.}

\noindent\textit{[From the Quarterly Journal of Pure and Applied Mathematics, vol.~xxi.\ (1886), pp.~92--107.]}

\medskip

\noindent\textbf{[Printed p.~344]}\hfill

In my Memoire ``sur les Hyperd\'eterminants,'' Crelle, t.~xxx.\ (1846), pp.~1--37, [13, 14, 16], I gave an investigation for the number of and relations between the quartic invariants of a given binary quantic: the very same formulae apply to the cubic seminvariants of a given binary quantic, or what is the same thing to the cubic seminvariants of a given weight, and I propose at present (considering the formulae in this point of view) to further develope the theory in regard to the solution of the systems of linear equations obtained in the memoir in question. But I first reproduce the investigation as it stands: I recall that the notation is the ordinary hyperdeterminant notation: for instance,
\[
\overline{12}^{2}\, UV = (\partial_{x_{1}}\partial_{y_{2}} - \partial_{x_{2}}\partial_{y_{1}})^{2}\, UV
\]
\[
\quad = (\partial_{x_{1}}^{2}\partial_{y_{2}}^{2} - 2\partial_{x_{1}}\partial_{y_{1}}\partial_{x_{2}}\partial_{y_{2}} + \partial_{x_{2}}^{2}\partial_{y_{1}}^{2})\, UV
\]
\[
\quad = a \cdot c - 2b \cdot b + c \cdot a
\]
\[
\quad = -2(ac - b^{2}),
\]
and that, when the variables do not disappear, each set $(x_{1}, y_{1})$, $(x_{2}, y_{2})$, \dots\ is ultimately replaced by $(x, y)$, so that these are the variables of any covariant.

The investigation* is as follows:

\footnotetext{[* This Collection, vol.~I., p.~104: see also, vol.~I., p.~117.]}

\medskip

``\dots We pass to the derivatives of the fourth degree, considering the forms in which all the differential coefficients are of the same order. It is easy to see that we may write
\[
\square\, UVWX = UVWX\,\overline{12}^{\alpha}\,\overline{13}^{\beta}\,\overline{14}^{\gamma}\,\overline{23}^{\beta'}\,\overline{24}^{\gamma'}\,\overline{34}^{\alpha'},
\]
\[
= D_{a, b, c}\, UVWX,\quad\text{or simply } D_{a, b, c}.
\]

\medskip

\noindent\textbf{[Printed p.~345]}\hfill

Writing for shortness
\[
\overline{12}\cdot\overline{34} = \mathfrak{A},\quad \overline{13}\cdot\overline{42} = \mathfrak{B},\quad \overline{14}\cdot\overline{23} = \mathfrak{C},
\]
we have
\[
D_{a, b, c} = \mathfrak{A}^{a}\mathfrak{B}^{b}\mathfrak{C}^{c}\, UVWX.
\]

Suppose $U = V = W = X$, and consider the derivatives which correspond to the same value of $\alpha$, $\beta$, $\gamma$: we have to find how many of these derivatives are independent, and to express the others in terms of these. Since the functions are equal after the differentiations, we may before the differentiations interchange in any manner the symbolic numbers $1$, $2$, $3$, $4$: this gives
\[
D_{a, b, c} = D_{b, a, c} = D_{a, c, b} = (-)^{f}\, D_{a, b, c} = (-)^{f}\, D_{b, a, c}.
\]

But we have identically
\[
\mathfrak{A} + \mathfrak{B} + \mathfrak{C} = 0.
\]

Multiplying by $\mathfrak{A}^{a}\mathfrak{B}^{b}\mathfrak{C}^{c}$, and applying it to the product $UVWX$, we have
\[
D_{a+1, b, c} + D_{a, b+1, c} + D_{a, b, c+1} = 0,
\]
which, putting $a + b + c = f - 1$, gives a system of equations between the derivatives $D_{\alpha\beta\gamma}$ for which $\alpha + \beta + \gamma = f$. Reducing these by the conditions just obtained, suppose that $\Theta f$ denotes the number of partitions of $f$ into three parts, zero included, and the permutations of the parts being disregarded, we have a number $\Theta f$ of derivatives $D_{\alpha\beta\gamma}$, and a number $\Theta(f - 1)$ of linear relations between these derivatives. There remains therefore a number $\Theta f - \Theta(f - 1)$ of independent derivatives: but when $f$ is an even number, we have included among these the functions $D_{\frac{1}{2}f, \frac{1}{2}f, 0}$, that is, $\overline{12}^{\frac{1}{2}f}\cdot\overline{34}^{\frac{1}{2}f}\, UVWX$, which evidently reduces itself to
\[
\overline{12}^{\frac{1}{2}f} UV \cdot \overline{34}^{\frac{1}{2}f} WX,\quad\text{or } B_{\frac{1}{2}f}(U, V)\cdot B_{\frac{1}{2}f}(W, X),
\]
that is to say, to the square of $B_{\frac{1}{2}f}(U, V)$. We must therefore diminish by unity this number $\Theta f - \Theta(f - 1)$, when $f$ is even. Let $E\!\left(\dfrac{f}{6}\right)$ denote the integer part of the fraction $\dfrac{f}{6}$: it can be shown that the required number is equal to
\[
E\!\left(\frac{f}{6}\right),\quad E\!\left(\frac{f + 3}{6}\right),
\]
according as $f$ is even or odd. Giving to $f$ the six forms
\[
6g,\ 6g + 1,\ 6g + 2,\ 6g + 3,\ 6g + 4,\ 6g + 5,
\]
we obtain for the independent derivatives of the fourth degree the corresponding numbers
\[
g,\ g,\ g,\ g + 1,\ g,\ g + 1;
\]
there is for example a single derivative for each of the orders $3$, $5$, $6$, $7$, $8$, $10$, two for each of the orders $9$, $11$, $12$, $13$, $14$, $16$, \&c.

\medskip

\noindent\textbf{[Printed p.~346]}\hfill

We may take for independent derivatives, when $f$ is even, the terms of the series $D_{f, 0, 0}$, $D_{f - 2, 2, 0}$, $D_{f - 4, 4, 0}$, \dots, and when $f$ is odd, those of the series $D_{f - 1, 1, 0}$, $D_{f - 3, 3, 0}$, $\dots$, continuing each series until the last term in which the first suffix exceeds or is equal to the second suffix. We have in each case the right number of terms; for example, in the case $f = 9$, we take for independent derivatives the two functions $D_{810}$ and $D_{630}$, and we form the equations
\begin{align*}
D_{900} + D_{810} + D_{801} &= 0, & D_{621} + D_{531} + D_{522} &= 0, \\
D_{810} + D_{720} + D_{711} &= 0, & D_{531} + D_{441} + D_{432} &= 0, \\
D_{801} + D_{711} + D_{702} &= 0, & D_{522} + D_{432} + D_{423} &= 0, \\
D_{720} + D_{630} + D_{621} &= 0, & D_{441} + D_{351} + D_{342} &= 0, \\
D_{711} + D_{621} + D_{612} &= 0, & D_{432} + D_{342} + D_{333} &= 0, \\
D_{702} + D_{612} + D_{603} &= 0, & &
\end{align*}
which are to be reduced by means of the formulas
\[
D_{900} = -D_{900} = 0,\quad D_{810} = -D_{810},\quad \&c.
\]

We will presently give the solution of these equations: but beginning with the second order and going successively to the ninth order, we form easily the following table:

\medskip

\begin{center}
\small
\begin{tabular}{lll}
$D_{200} = B^{2}$, & $D_{500} = 0$, & $D_{700} = 0$, \\
$D_{110} = -\tfrac{1}{2}B^{2}$, & $D_{410} = B^{2}$, & $D_{610} = B^{2}$, \\
& $D_{320} = -\tfrac{1}{2}B^{2}$, & $D_{520} = -\tfrac{1}{2}B^{2}$, \\
& $D_{311} = 0$, & $D_{511} = 0$, \\
$D_{300} = 0$, & $D_{221} = 0$, & $D_{430} = D_{610} - \tfrac{1}{2}D_{540}$, \\
$D_{210} = 0$, & & $D_{421} = -\tfrac{1}{2}D_{610} - \tfrac{1}{2}D_{540}$, \\
$D_{111} = 0$, & & $D_{331} = -\tfrac{1}{2}D_{540}$, \\
& & $D_{322} = -\tfrac{1}{2}D_{610} - \tfrac{1}{4}D_{540}$, \\
$D_{400} = B^{2}$, & $D_{600} = B^{2}$, & \\
$D_{310} = -\tfrac{1}{2}B^{2}$, & $D_{510} = -\tfrac{1}{2}B^{2}$, & \\
$D_{220} = \tfrac{1}{3}B^{2}$, & $D_{420} = \tfrac{1}{3}B^{2}$, & \\
$D_{211} = 0$, & $D_{411} = -\tfrac{1}{6}B^{2}$, & \\
& $D_{330} = -\tfrac{1}{3}B^{2}$, & \\
& $D_{321} = \tfrac{1}{6}B^{2}$, & \\
& $D_{222} = -\tfrac{1}{3}B^{2}$, & \\
\end{tabular}
\end{center}

\medskip

\noindent\textbf{[Printed p.~347]}\hfill

For any value of $f$ except $f = 2$, $3$, or $4$, the table commences, for $f$ even and $f$ odd respectively, in the following manner:
\[
D_{f, 0, 0} = B_{\frac{1}{2}f}^{2},\qquad D_{f, 0, 0} = 0,
\]
\[
D_{f - 1, 1, 0} = -\tfrac{1}{2}B_{\frac{1}{2}f}^{2},\qquad D_{f - 1, 1, 0} = 0,
\]
\[
D_{f - 2, 2, 0} + D_{f - 2, 1, 1} = D_{f - 2, 2, 0} = -D_{f - 1, 1, 0},
\]
\[
D_{f - 3, 3, 0} = D_{f - 2, 1, 1} = 0,
\]
but beyond this I do not know the law of the series.''

Before going further, I remark that if $\mathfrak{A}$, $\mathfrak{B}$, $\mathfrak{C}$, instead of the foregoing values, denote respectively
\begin{align*}
\mathfrak{A} &= (\alpha\partial_{x} + \beta\partial_{y})\,\overline{23},\\
\mathfrak{B} &= (\alpha\partial_{x_{2}} + \beta\partial_{y_{2}})\,\overline{31},\\
\mathfrak{C} &= (\alpha\partial_{x_{3}} + \beta\partial_{y_{3}})\,\overline{12},
\end{align*}
then identically
\[
\mathfrak{A} + \mathfrak{B} + \mathfrak{C} = 0,
\]
and the same theory applies to the cubic derivatives $\mathfrak{A}^{a}\mathfrak{B}^{b}\mathfrak{C}^{c}\, UVW$, that is, to the cubic covariants, or, attending only to the coefficients of the highest powers of $x$, to the cubic seminvariants.

Or, we may use the Clebschian notation, for instance
\[
(a, b, c\,|\,x, y)^{2} = (a_{0}x + a_{1}y)^{2} = (b_{0}x + b_{1}y)^{2} = \text{\&c.},
\]
that is,
\[
a_{0}^{2} = b_{0}^{2} = \dots = a,
\]
\[
a_{0}a_{1} = b_{0}b_{1} = \dots = b,
\]
\[
a_{1}^{2} = b_{1}^{2} = \dots = c,
\]
viz.\ in the language of Prof.\ Sylvester, $a_{0}$, $a_{1}$, $b_{0}$, $b_{1}$, \dots\ are here \textit{umbrae}. The invariant is
\[
(a_{0}b_{1} - a_{1}b_{0})^{2} = a_{0}^{2}b_{1}^{2} - 2a_{0}a_{1}b_{0}b_{1} + a_{1}^{2}b_{0}^{2},
\]
\[
= ac - 2b\cdot b + ac,
\]
\[
= -2(ac - b^{2}),
\]
and so in other cases. To apply this directly to the theory of seminvariants, it is more convenient to write
\[
(1, b, c\,|\,x, y)^{2} = (x + \alpha y)^{2} = (x + \beta y)^{2} = \&c.,
\]
that is,
\[
\alpha = \beta = \dots = b,
\]
\[
\alpha^{2} = \beta^{2} = \dots = c,
\]
where $\alpha$, $\beta$, \dots\ are \textit{umbrae}.

The invariant is
\[
(\alpha - \beta)^{2} = \alpha^{2} - 2\alpha\beta + \beta^{2}
\]
\[
= c - 2b\cdot b + c
\]
\[
= 2(c - b^{2});
\]

\medskip

\noindent\textbf{[Printed p.~348]}\hfill

\noindent or, to take the case of a cubic function
\[
(1, b, c, d\,|\,x, y)^{3} = (x + \alpha y)^{3} = (x + \beta y)^{3} = (x + \gamma y)^{3},
\]
that is,
\[
\alpha = \beta = \gamma = b,
\]
\[
\alpha^{2} = \beta^{2} = \gamma^{2} = c,
\]
\[
\alpha^{3} = \beta^{3} = \gamma^{3} = d,
\]
a seminvariant is
\begin{align*}
(\alpha - \beta)^{2}(\alpha - \gamma) &= \alpha^{3} - 2\alpha^{2}\beta + \alpha\beta^{2} - \alpha^{2}\gamma + 2\alpha\beta\gamma - \beta^{2}\gamma \\
&= d - 2bc + bc - bc + 2b^{3} - bc \\
&= d - 3bc + 2b^{3},
\end{align*}
and so in other cases. Writing here
\[
\mathfrak{A} = \beta - \gamma,\quad \mathfrak{B} = \gamma - \alpha,\quad \mathfrak{C} = \alpha - \beta,
\]
the general form of a cubic seminvariant of the weight $\omega$ is now $\mathfrak{A}^{p}\mathfrak{B}^{q}\mathfrak{C}^{r}$, or say $D_{p, q, r}$, where $p + q + r = \omega$ (the new letters $p$, $q$, $r$ being used for the exponents, since $\alpha$, $\beta$, $\gamma$ are now employed in a different sense; and since $f$ will be required as a coefficient, I have also written $\omega$ instead of $f$ for the weight). We have, as before,
\[
\mathfrak{A} + \mathfrak{B} + \mathfrak{C} = 0,
\]
and we again see that the theory as to the number of and relations between the cubic seminvariants is identical with that of the quartic invariants. Observe also that $D_{\omega, 0, 0} = 0$, for $\omega$ odd, while $D_{\omega, 0, 0} = \mathfrak{A}^{\omega}$, for $\omega$ even, is a quadric seminvariant, which is of course not to be reckoned among the proper cubic seminvariants; this exactly corresponds to the $B_{\frac{1}{2}f}^{2}$ which is not a proper quartic invariant.

The choice of the functions $D_{f, 0, 0}$, $D_{f - 2, 2, 0}$, \dots, or $D_{f - 1, 1, 0}$, $D_{f - 3, 3, 0}$, \dots, for expressing in terms of them the other functions has its advantages, but also its disadvantages; in what follows, I in effect replace them by $D_{f}$, $D_{f - 2, 2, 0}$, $D_{f - 4, 4, 0}$, and $D_{f - 1, 1, 0}$, $D_{f - 3, 3, 0}$, \dots\ respectively. And instead of considering the whole series of the functions $D_{p, q, r}$ for a given weight $\omega$, I consider only those of the form $D_{p, q, 0}$; these form a single series $D_{\omega, 0, 0}$, $D_{\omega - 1, 1, 0}$, $D_{\omega - 2, 2, 0}$, \dots; and for shortness, I call them $A$, $B$, $C$, $D$, \&c.\ respectively, viz.\ I write
\[
A,\ B,\ C,\ D,\ \dots = (\mathfrak{A}^{\omega},\ \mathfrak{A}^{\omega - 1}\mathfrak{B},\ \mathfrak{A}^{\omega - 2}\mathfrak{B}^{2},\ \mathfrak{A}^{\omega - 3}\mathfrak{B}^{3},\ \dots).
\]

Suppose first that $\omega$ is odd; we have
\[
D_{\omega - q - r, q, r} = (-)^{r}\, D_{\omega - q - r, r, q},
\]
and, in particular,
\[
D_{\omega - 2q, q, q} = -D_{\omega - 2q, q, q},
\]
that is,
\[
D_{\omega - 2q, q, q} = 0,
\]
that is,
\[
\mathfrak{A}^{\omega - 2q}\,\mathfrak{B}^{q}\mathfrak{C}^{q} = 0;
\]
Hence
\[
A = 0;
\]

\medskip

\noindent\textbf{[Printed p.~349]}\hfill

\noindent moreover
\[
B + C = \mathfrak{A}^{\omega - 1}\mathfrak{B} + \mathfrak{A}^{\omega - 2}\mathfrak{B}^{2} = \mathfrak{A}^{\omega - 2}\mathfrak{B}(\mathfrak{A} + \mathfrak{B}) = -\mathfrak{A}^{\omega - 2}\mathfrak{B}\mathfrak{C} = 0;
\]
and similarly $C + 2D + E = 0$, \&c.; that is, we have
\begin{align*}
A &= 0, \\
B + C &= 0, \\
C + 2D + E &= 0, \\
D + 3E + 3F + G &= 0, \\
&\text{\&c.}
\end{align*}

The readiest way of solving these is to express the other functions in terms of $B$, $D$, $F$, $H$, \&c.; viz.\ we thus have

\medskip

\begin{center}
\small
\begin{tabular}{|c|cccccccccccc|}
\hline
 & $A$ & $B$ & $C$ & $D$ & $E$ & $F$ & $G$ & $H$ & $I$ & $J$ & $K$ & $L$ \\
\hline
$B$ & $0$ & $1$ & $-1$ & $0$ & $1$ & $0$ & $-3$ & $0$ & $17$ & $0$ & $-155$ & $0$ \\
$D$ & & & & $1$ & $-2$ & $0$ & $5$ & $0$ & $-28$ & $0$ & $255$ & $0$ \\
$F$ & & & & & & $1$ & $-3$ & $0$ & $14$ & $0$ & $-126$ & $0$ \\
$H$ & & & & & & & & $1$ & $-4$ & $0$ & $30$ & $0$ \\
$J$ & & & & & & & & & & $1$ & $-5$ & $0$ \\
$L$ & & & & & & & & & & & & $1$ \\
\hline
\end{tabular}
\end{center}

\medskip

\noindent read according to the columns
\begin{align*}
A &= 0, \\
B &= B, \\
C &= -B, \\
D &= D, \\
E &= B - 2D, \\
F &= F, \\
G &= -3B + 5D - 3F, \\
&\text{\&c.}
\end{align*}

\medskip

\noindent and, by way of verification of the numbers, observe that the sum of the numbers in a column is $1$ and $-1$ alternately.

The formulae are true for any odd value of $\omega$ whatever; but they require an explanation, viz.\ for any finite value of $\omega$, the terms $B$, $D$, $F$, \dots\ are not all of them arbitrary. For any given value of $\omega$ the number of arbitrary terms is in fact given by what precedes, and it is also known by Captain MacMahon's theorem, viz.\ the number of cubic seminvariants is = that of the non-unitary symmetric functions $3^{a}2^{b}$ of the proper weight $3a + 2b = \omega$; thus for $\omega = 9$, the forms are $3222$ and $333$, so that there should be only two arbitrary terms.

\medskip

\noindent\textbf{[Printed p.~350]}\hfill

We have here, writing for shortness $900$, $810$, \&c., instead of $D_{900}$, $D_{810}$, \&c.\ respectively,
\begin{align*}
900 &= 0, \\
810 &= B, \\
720 &= -B, \\
630 &= D, \\
540 &= B - 2D, \\
450 &= F, \\
360 &= -3B + 5D - 3F, \\
270 &= H, \\
180 &= 17B - 28D + 14F - 4H, \\
090 &= L.
\end{align*}

But we have $900 = 0$, $810 + 180 = 0$, $720 + 270 = 0$, $630 + 360 = 0$, $540 + 450 = 0$; that is,
\begin{align*}
L &= 0, \\
18B - 28D + 14F - 4H &= 0, \\
-B + H &= 0, \\
-3B + 6D - 3F &= 0, \\
B - 2D + F &= 0,
\end{align*}
all satisfied by $F = B + 2D$, $H = B$, $L = 0$. The proper course is to stop at the equation for $540$, viz.\ the system is
\begin{align*}
900 &= 0, \\
810 &= B, \\
720 &= -B, \\
630 &= D, \\
540 &= B - 2D,
\end{align*}
equations which may be considered as expressing the several functions $900$, $810$, $720$, $630$, $540$ in terms of $B = 810$ and $D = 630$. They agree, as they should do, with a foregoing result,
\begin{align*}
900 &= 0, \\
810 &= 810, \\
720 &= -810, \\
630 &= \tfrac{1}{2}\cdot 810 - \tfrac{1}{2}\cdot 540, \\
540 &= 540,
\end{align*}
and it would of course be possible to express the remaining forms $711$, $621$, $522$, $441$, $432$, and $333$ in terms of $B = 810$ and $D = 630$. But observe that the speciality of

\medskip

\noindent\textbf{[Printed p.~351]}\hfill

\noindent the present process is that, instead of the whole series of forms $900$, $810$, $720$, $711$, $630$, $621$, $540$, $531$, $522$, $441$, $432$, $330$, we work only with the forms $900$, $810$, $720$, $630$, $540$, for which the last element is $= 0$.

A more simple solution of the system of linear equations in $A$, $B$, $C, \dots$ is the following:

\medskip

\begin{center}
\small
\begin{tabular}{|c|cccccccccccc|}
\hline
 & $A$ & $B$ & $C$ & $D$ & $E$ & $F$ & $G$ & $H$ & $I$ & $J$ & $K$ & $L$ \\
\hline
$x$ & $0$ & $1$ & $-1$ & $-1$ & $1$ & $-1$ & $1$ & $-1$ & $1$ & $-1$ & $1$ & $-1$ \\
$y$ & & & & $1$ & $-2$ & $3$ & $-4$ & $5$ & $-6$ & $7$ & $-8$ & $9$ \\
$z$ & & & & & & $1$ & $-3$ & $5$ & $-10$ & $15$ & $-21$ & $28$ \\
$u$ & & & & & & & & $1$ & $-4$ & $10$ & $-20$ & $35$ \\
$v$ & & & & & & & & & & $1$ & $-5$ & $15$ \\
$w$ & & & & & & & & & & & & $1$ \\
\hline
\end{tabular}
\end{center}

\medskip

\noindent read
\begin{align*}
A &= 0, \\
B &= x, \\
C &= -x, \\
D &= x + y, \\
E &= -x - 2y, \\
F &= x + 3y + z, \\
G &= -x - 4y - 3z,
\end{align*}
where $x$, $y$, $z, \dots$ are arbitrary; the numbers in the table are the binomial coefficients with the signs $+$ and $-$ alternately.

We may, it is clear, express $x$, $y$, $z, \dots$ in terms of $B$, $D$, $F, \dots$, viz.\ we thus have
\begin{align*}
x &= B, \\
y &= -B + D, \\
z &= 2B - 3D + F,
\end{align*}
and substituting these values, we reproduce the foregoing expressions of $A$, $B$, $C$, \dots\ in terms of $B$, $D$, $F$, \dots.

For a given finite value of $\omega$, we have of course the same number of arbitrary coefficients $x$, $y$, $z$, \dots\ as there were of arbitrary coefficients $B$, $D$, $F$, \dots; thus for $\omega = 9$, only $x$ and $y$ are arbitrary, and the remaining coefficients $z$, \&c., are determined in terms of these. Or, what is the same thing, we stop with the equation $E = -x - 2y$, next preceding an equation with the non-arbitrary coefficient $z$.

\medskip

\noindent\textbf{[Printed p.~352]}\hfill

For the case $\omega$ even, I write
\begin{align*}
A + 2B &= A', \\
B + 2C &= B', \\
C + 2D &= C', \\
&\text{\&c.;}
\end{align*}
and I say that the new functions $A'$, $B'$, $C', \dots$, are related together precisely in the same manner as are the functions $A$, $B$, $C, \dots$, belonging to an odd weight $\omega$; viz.\ we have
\begin{align*}
A' &= 0, \\
B' + C' &= 0, \\
C' + 2D' + E' &= 0, \\
&\text{\&c.,}
\end{align*}
which being so, the theory is included in what precedes, and there is no occasion to consider it separately.

To prove this, observe that $\omega$ being even, we have
\[
D_{\omega - q - r, q, r} = D_{\omega - q - r, r, q},
\]
that is,
\[
\mathfrak{A}^{\omega - q - r}\,\mathfrak{B}^{q}\mathfrak{C}^{r} = \mathfrak{A}^{\omega - q - r}\,\mathfrak{B}^{r}\mathfrak{C}^{q},
\]
whence, in particular,
\[
\mathfrak{C}^{q} = \mathfrak{B}^{q}\quad (\omega - 2q\ \text{even});
\]
we thus have
\[
\mathfrak{C} = \mathfrak{B},\quad \mathfrak{C}^{2} = \mathfrak{A}^{\omega - 1}\mathfrak{B} = \mathfrak{A}^{\omega - 2}\mathfrak{B}^{2},\quad \mathfrak{C}^{3} = \mathfrak{A}^{\omega - 3}\mathfrak{B}^{3},\ \&\text{c.};
\]
\[
A' = A + 2B = \mathfrak{A}^{\omega} + 2\mathfrak{A}^{\omega - 1}\mathfrak{B} = \mathfrak{A}^{\omega} + \mathfrak{A}^{\omega - 1}\mathfrak{B} + \mathfrak{A}^{\omega - 1}\mathfrak{C} = \mathfrak{A}^{\omega - 1}(\mathfrak{A} + \mathfrak{B} + \mathfrak{C}),
\]
which is $= 0$. Similarly
\[
B' + C' = B + 3C + 2D = \mathfrak{A}^{\omega - 1}\mathfrak{B} + 3\mathfrak{A}^{\omega - 2}\mathfrak{B}^{2} + 2\mathfrak{A}^{\omega - 3}\mathfrak{B}^{3},
\]
or, since
\[
B + C = \mathfrak{A}^{\omega - 2}\mathfrak{B}(\mathfrak{A} + 3\mathfrak{B} + 2\mathfrak{B}^{2}) = \mathfrak{A}^{\omega - 2}\mathfrak{B}(\mathfrak{A} + 2\mathfrak{B}),
\]
this is
\[
= -\mathfrak{A}^{\omega - 2}\mathfrak{B}\mathfrak{C}(\mathfrak{A} + 2\mathfrak{B}),
\]
which, in virtue of
\[
\mathfrak{A} + \mathfrak{B} + \mathfrak{C} = 0,
\]
is
\[
= -\mathfrak{A}^{\omega - 2}\mathfrak{B}\mathfrak{C}(-\mathfrak{C}) = \mathfrak{A}^{\omega - 2}\mathfrak{B}\mathfrak{C}^{2},
\]
which is $= 0$. And similarly $C' + 2D' + E' = 0$, \&c.

I come to a new part of the theory: to fix the ideas, consider the weight $\omega = 27$; and write also
\[
a_{0},\ a_{1},\ a_{2},\ a_{3}, \dots = 1,\ b,\ c,\ d,\ e,\ f,\ g,\ h,\ i,\ j,\ k,\ l,\ m, \dots,
\]

\medskip

\noindent\textbf{[Printed p.~353]}\hfill

\noindent using, if we please, the suffixed $a$ for the higher terms. The cubic seminvariants are $3^{9} \cdot 2^{0}$, $3^{7}\cdot 2^{3}$, $3^{5}\cdot 2^{6}$, $3^{3}\cdot 2^{9}$, $3\cdot 2^{12}$, viz.\ the number of them is $5$; we have thus the $5$ terms $B$, $D$, $F$, $H$, $J$. Here
\[
B = (\alpha - \beta)^{26}(\alpha - \gamma),
\]
or, interchanging the $\alpha$ and $\beta$,
\[
= (\beta - \alpha)^{26}(\beta - \gamma) = -(\alpha - \beta)^{26}(\beta - \gamma);
\]
there is an initial term $\alpha^{27}$ and a final term $\beta^{26}\gamma$, $= b^{26}c$ or $bn_{2}$, or we may write
\[
B = (\alpha - \beta)^{26}(\alpha - \gamma) = a_{27} + bn_{2};
\]
and similarly for the other forms. We have thus
\begin{align*}
B &= (\alpha - \beta)^{26}(\alpha - \gamma) = a_{27} + bn_{2}, \\
D &= (\alpha - \beta)^{24}(\alpha - \gamma)^{3} = a_{27} + dm_{2}, \\
F &= (\alpha - \beta)^{22}(\alpha - \gamma)^{5} = a_{27} + fl_{2}, \\
H &= (\alpha - \beta)^{20}(\alpha - \gamma)^{7} = a_{27} + hk_{2}, \\
J &= (\alpha - \beta)^{18}(\alpha - \gamma)^{9} = a_{27} + j^{3},
\end{align*}
where the initial term $a^{27} = a_{27}$ has in each case the coefficient unity, but the final terms have each of them the proper numerical coefficient.

It must be possible to form linear combinations
\begin{align*}
B &= a_{27} + bn_{2}, \\
(B, D) &= a_{25}(c - b^{2}) + dm_{2}^{3}, \\
(B, D, F) &= a_{25}(e + c^{2}) + fl_{2}, \\
(B, D, F, H) &= a_{21}(g + d^{2}) + hk^{3}, \\
(B, D, F, H, J) &= a_{19}(i + e^{2}) + j^{3},
\end{align*}
where $c - b^{2}$, $e + c^{2}$, $g + d^{2}$, $i + e^{2}$ denote the seminvariants
\[
c - b^{2},\quad e - 4bd + 3c^{2},\quad g - 6bf + 15ce - 10d^{2},
\]
\[
i - 8bh + 28cg - 56df + 35e^{2},
\]
respectively. The expressions indicated by their initial and final terms $a_{27} + bn_{2}$, $a_{25}c + dm_{2}^{3}$, \&c., are what I call columns;\textsuperscript{*} see my paper ``Seminvariant Tables,'' \textit{American Journal of Mathematics}, t.~VII.\ (1885), pp.~59--73, [831], where, however, the theory is not by any means completely developed.

As to this, observe that $B$, $D$ are each of them of the form $a_{27} + a_{0}^{0} a_{25}(c, b^{2})\dots$, the proper combination in order to eliminate $a_{27}$ is obviously $B - D$, the term in $a_{0}b$ will then disappear of itself, and the terms in $a_{25}(c, b^{2})$ will combine into a term $a_{25}(c - b^{2})$, viz.\ in the function qua seminvariant the coefficient of the highest letter $a_{25}$ will be the seminvariant $c - b^{2}$. Similarly, in the combination $B$, $D$, $F$, the two coefficients will be determinable so that there are no terms in $a_{0}$, $a_{1}$, $a_{25}$ or $a_{0}d$, and this being so, the term in $a_{23}$ will be $a_{23}(e + 4bd - 3c^{2})$, viz.\ in the function, qua seminvariant, the coefficient of the highest letter $a_{23}$ will be the seminvariant $e - \dots$. And similarly for the remaining two linear combinations.

\medskip

\noindent\textbf{[Printed p.~354]}\hfill

\noindent As a partial verification, I write
\[
B = (\alpha - \beta)^{26}(\alpha - \gamma)
\]
\[
\quad = \{\alpha^{26} + \beta^{26} - 26(\alpha^{25}\beta + \alpha\beta^{25}) + 325(\alpha^{24}\beta^{2} + \alpha^{2}\beta^{24})\dots\}(\alpha - \gamma),
\]
or, retaining only the terms which have an index at least $= 25$, this is
\begin{align*}
B &= \alpha^{27} + \alpha\beta^{26} - 26\alpha^{26}\beta + 325\alpha^{25}\beta^{2} \\
   &\quad - (\alpha^{26} + \alpha^{25}\beta)\gamma \\
   &\quad = a_{27} + a_{0}b \\
   &\quad\quad - 26a_{0}b - 26a_{25}c \\
   &\quad\quad - 26a_{0}b + 52a_{25}b^{2} \\
   &= a_{27} - 27a_{0}b + a_{25}(299c - 52b^{2});
\end{align*}
and similarly
\[
D = (\alpha - \beta)^{24}(\alpha - \gamma)^{3}
\]
\[
\quad = \{\alpha^{24} + \beta^{24} - 24(\alpha^{23}\beta + \alpha\beta^{23}) + 276(\alpha^{22}\beta^{2} + \alpha^{2}\beta^{22})\dots\}(\alpha - \gamma)^{3},
\]
or, retaining only the terms which have an index at least $= 25$, this is
\begin{align*}
D &= \alpha^{27} - 24\alpha^{26}\beta + 276\alpha^{25}\beta^{2} \\
   &\quad - 3\alpha^{26}\gamma + 72\alpha^{25}\beta\gamma \\
   &\quad + 3\alpha^{25}\gamma^{2} \\
   &\quad = a_{27} - 24a_{0}b + 276a_{25}b^{2} \\
   &\quad\quad - 3a_{0}b + 72a_{25}b^{2} \\
   &\quad\quad + 3a_{25}c \\
   &\quad = a_{27} + a_{25}(279c + 72b^{2}),
\end{align*}
and thence
\[
B - D = a_{25}(20c - 20b^{2}),
\]
viz.\ the coefficient of $a_{25}$ is $= 20(c - b^{2})$.

But instead of the foregoing forms $B$, $D$, $F$, $H$, $J$, we may start with five other forms which are in fact linear combinations of these, viz.\ these are the forms
\begin{align*}
X &= (\alpha - \beta)^{26}(\alpha + \beta - 2\gamma) = a_{27} + bn_{2}, \\
Y &= (\alpha - \beta)^{24}(\alpha + \beta - 2\gamma)(\alpha - \gamma)(\beta - \gamma) = a_{25}(c - b^{2}) + dm_{2}, \\
Z &= (\alpha - \beta)^{22}(\alpha + \beta - 2\gamma)(\alpha - \gamma)^{2}(\beta - \gamma)^{2} = a_{25}(c - b^{2}) + fl_{2}, \\
W &= (\alpha - \beta)^{20}(\alpha + \beta - 2\gamma)(\alpha - \gamma)^{3}(\beta - \gamma)^{3} = a_{23}(e + c^{2}) + hk_{2}, \\
T &= (\alpha - \beta)^{18}(\alpha + \beta - 2\gamma)(\alpha - \gamma)^{4}(\beta - \gamma)^{4} = a_{23}(e^{2} + c^{2}) + j^{3},
\end{align*}
where as well the initial as the final terms should have their proper numerical coefficients.

\medskip

\noindent\textbf{[Printed p.~355]}\hfill

As to this, observe that in $Y$ we have a term $\alpha^{26}(\beta - \gamma)$ which is $= 0$, and similarly a term $\beta^{26}(\alpha - \gamma)$ which is $= 0$; the highest powers are thus $\alpha^{25}$ and $\beta^{25}$, giving a term in $a_{25}$ which can only be $a_{25}(c - b^{2})$. In $Z$ we have the terms $\alpha^{25}(\beta - \gamma)^{2}$ and $\beta^{25}(\alpha - \gamma)^{2}$, giving the term $a_{25}(c - b^{2})$; in $W$ we have the terms $\alpha^{24}(\beta - \gamma)^{3}$ and $\beta^{24}(\alpha - \gamma)^{3}$ which are each $= 0$, the highest powers are thus $\alpha^{23}$, $\beta^{23}$, giving a term in $a_{23}$ which can only be $a_{23}(e + 4bd - 3c^{2})$; and in $T$ we have the terms $\alpha^{4}(\beta - \gamma)^{4}$ and $\beta^{4}(\alpha - \gamma)^{4}$ which give the term $a_{23}(e + 4bd - 3c^{2})$.

The combinations which give the columns thus are
\begin{align*}
X &= a_{27} + bn_{2}, \\
(Y) &= a_{25}(c - b^{2}) + dm_{2}, \\
(Y, Z) &= a_{23}(e + c^{2}) + fl_{2}, \\
(Y, Z, W) &= a_{21}(g + d^{2}) + hk_{2}, \\
(Y, Z, W, T) &= a_{19}(i + e^{2}) + j^{3}.
\end{align*}
I start with the form
\[
u = (Y,\ Z,\ W,\ T) = uY + \lambda Z + \mu W + \nu T.
\]

We know that it is possible to determine the coefficients $\lambda$, $\mu$, $\nu$, in suchwise that the linear function shall be of the proper form $a_{19}(i + e^{2}) + j^{3}$; or, what is the same thing, that there shall be no term with an index exceeding $19$; the conditions to be satisfied are apparently more than three, but of course the number of independent conditions must be $= 3$. We have, in particular, to get rid of all the terms of the second degree which precede $a_{19}i$, viz.\ the coefficients of $a_{19}c$, $a_{24}d$, $a_{23}e$, $a_{22}f$, $a_{21}g$, $a_{20}h$ must all of them vanish. Now the terms of $u$ which produce terms of the second degree are the terms containing any two but not all three of the symbolic letters $\alpha$, $\beta$, $\gamma$, say the biliteral terms; we have for instance $\alpha^{25}\beta^{2}$, $\alpha^{25}\gamma^{2}$, $\beta^{25}\gamma^{2}$ (there are no terms $\gamma^{25}\alpha^{2}$ or $\gamma^{25}\beta^{2}$, since the highest power of $\gamma$ is $\gamma^{9}$), each of which is $a_{25}c$. But in the development of $u$, we may in any biliteral term by an interchange of the letters $\alpha$, $\beta$, $\gamma$ make the letter of highest index to be $\alpha$, and the other letter to be $\beta$; thus the several terms $\alpha^{25}\gamma^{2}$, $\beta^{25}\alpha^{2}$, $\beta^{25}\gamma^{2}$ may be each of them converted into $\alpha^{25}\beta^{2}$, and so in other cases. Imagining this to be done, the term in $a_{25}c$ is simply the term in $\alpha^{25}\beta^{2}$, and in like manner the term in $a_{24}d$ is the term in $\alpha^{24}\beta^{3}$, and so in other cases; the condition as to the disappearance of the terms $a_{25}c$, \dots, $a_{20}h$ is the condition that the terms in $\alpha^{25}\beta^{2}$, $\alpha^{24}\beta^{3}$, $\alpha^{23}\beta^{4}$, $\alpha^{22}\beta^{5}$, $\alpha^{21}\beta^{6}$, $\alpha^{20}\beta^{7}$ shall all of them disappear. And if, in the function $u$ transformed in the foregoing manner, we write for convenience $\alpha = 1$, so that $u$ has become a function of $\beta$ only (and observe that $u$ will contain the factor $\beta^{2}$), then the condition is that the terms in $\beta^{2}$, $\beta^{3}$, $\beta^{4}$, $\beta^{5}$, $\beta^{6}$, $\beta^{7}$ may in fact be written $u = 0$; viz.\ it is assumed that $u$ is in the first instance transformed into a function of $\beta$ as just explained, and the equation is then to be understood as denoting that the coefficients of $u$ are to be so determined that as many as possible of the terms (beginning with that which contains the lowest power of $\beta$) shall each of them vanish.

Consider any term of $u$, for instance the term
\[
F = (\alpha - \beta)^{24}(\alpha + \beta - 2\gamma)(\alpha - \gamma)(\beta - \gamma).
\]

\end{document}
